This diploma thesis presents the redesign and extension of the \emph{Tumbler~v3.0} mobile robot platform at HTL~BULME Graz-Gösting into an autonomous testbed for self-parking experiments. To eliminate wiring faults and mechanical fragility, the previous modular breadboard-style electronics are replaced by a custom four-layer mainboard centered around an \emph{ESP32-S3-WROOM-1} microcontroller, integrating dual \emph{TB67H450FNG} H-bridge motor drivers, an \emph{LSM6DSV} six-axis inertial measurement unit, power-protection circuitry, and USB-C connectivity, alongside a redesigned six-channel \emph{CNY70} optical line-sensor board. Environmental perception is achieved by an array of eight \emph{VL53L0/1XV2} time-of-flight LiDAR sensors embedded into a newly engineered, crash-resistant unibody chassis (Revision~v9.0) verified for zero geometric interference across all 121~CAD bodies. On the software side, raw distance readings are filtered in real time to extract candidate parking spaces by width, depth, and orientation, feeding a state machine and closed-loop motor controller to execute collision-free parking maneuvers.
